%% Reference deck — every supported slide pattern, once.
%% Copy a frame from here rather than writing one from memory.
\documentclass[aspectratio=169]{beamer}

\usepackage{Archivo}
\renewcommand{\sfdefault}{\rmdefault}
\usepackage{graphicx}
\usepackage{tikz}
\usetikzlibrary{calc,positioning}
\tikzset{
  invisible/.style={opacity=0, text opacity=0},
  visible on/.style={alt={#1{}{invisible}}},
  alt/.code args={<#1>#2#3}{%
    \alt<#1>{\pgfkeysalso{#2}}{\pgfkeysalso{#3}}%
  },
}
\usepackage[most]{tcolorbox}
\usepackage{listings}

%% ---- Superteam Poland brand colors ----
\definecolor{stpldark}{HTML}{0B0E0F}   % near-black background
\definecolor{stplcard}{HTML}{161A1B}   % card / block surface
\definecolor{stplborder}{HTML}{1E2526}  % subtle track color for progress bar
\definecolor{stplred}{HTML}{DC143C}    % brand crimson red
\definecolor{stplwhite}{HTML}{FFFFFF}  % white
\definecolor{stplmuted}{HTML}{8A9BA8}  % muted / secondary text

\usetheme{default}
\setbeamersize{text margin left=0.85cm, text margin right=0.85cm}

%% ---- Background & base text ----
\setbeamercolor{background canvas}{bg=stpldark}
\setbeamercolor{normal text}{fg=stplwhite, bg=stpldark}
\usebeamercolor{normal text}

%% ---- Background image helpers ----
\newcommand{\stpldefaultbg}{%
  \setbeamertemplate{background}{%
    \begin{tikzpicture}[remember picture, overlay]
      \node[at=(current page.center), inner sep=0pt] {%
        \includegraphics[width=\paperwidth,height=\paperheight]{images/bg-default.png}%
      };
      \fill[stplborder]
        (current page.north west) rectangle
        ([yshift=-2.5pt]current page.north east);
      \pgfmathsetmacro{\stplprog}{\insertframenumber/\inserttotalframenumber}
      \fill[stplred]
        (current page.north west) rectangle
        ($(current page.north west)+(\stplprog\paperwidth,-2.5pt)$);
    \end{tikzpicture}%
  }%
}
\newcommand{\stplherobg}{%
  \setbeamertemplate{background}{%
    \includegraphics[width=\paperwidth,height=\paperheight]{images/bg-hero.png}%
  }%
}
\stpldefaultbg

%% ---- Structure ----
\setbeamercolor{structure}{fg=stplred}

%% ---- Title slide metadata ----
\setbeamercolor{title}{fg=stplwhite}
\setbeamercolor{subtitle}{fg=stplred}
\setbeamercolor{author}{fg=stplwhite}
\setbeamercolor{institute}{fg=stplmuted}
\setbeamercolor{date}{fg=stplmuted}

%% ---- Frame title: flat, no box background ----
\setbeamercolor{frametitle}{fg=stplwhite, bg=}
\setbeamerfont{frametitle}{size=\large, series=\bfseries}
\setbeamertemplate{frametitle}{%
  \vspace{0.35cm}%
  {\centering\insertframetitle\par}%
  \nointerlineskip\vspace{0.18cm}%
  \hfil{\color{stplred}\rule{2.2cm}{1.5pt}}\hfil\par%
}



%% ---- Incremental reveals ----
%% Bullets reveal one at a time by default. See style.md §4.5 for the two
%% per-frame opt-outs (static frame: `\begin{frame}<1>{...}`; manually
%% choreographed frame: `\beamerdefaultoverlayspecification{}` at frame top).
\beamerdefaultoverlayspecification{<+->}

%% ---- Bullet items ----
\setbeamercolor{itemize item}{fg=stplred}
\setbeamercolor{itemize subitem}{fg=stplmuted}
\setbeamercolor{itemize subsubitem}{fg=stplmuted}
\setbeamertemplate{itemize item}{\bfseries\textendash}
\setbeamertemplate{itemize subitem}{\small\textendash}
\setbeamercolor{enumerate item}{fg=stplred}
\setbeamercolor{enumerate subitem}{fg=stplmuted}
\setbeamertemplate{enumerate item}{\textcolor{stplred}{\bfseries\insertenumlabel.}}
\setbeamertemplate{enumerate subitem}{\textcolor{stplmuted}{\small\insertenumlabel.}}

%% ---- Blocks (flat left-accent via tcolorbox) ----
\tcbset{
  stplbase/.style={
    enhanced, breakable,
    colback=stplcard, coltext=stplwhite,
    coltitle=stplwhite, fonttitle=\bfseries,
    boxrule=0pt, leftrule=3pt,
    arc=0pt, outer arc=0pt,
    top=5pt, bottom=5pt, left=7pt, right=7pt,
    toptitle=4pt, bottomtitle=4pt,
    titlerule=0pt,
  },
}
\renewenvironment{block}[1]{%
  \pause%
  \begin{tcolorbox}[stplbase, colframe=stplred, title={#1}]%
}{%
  \end{tcolorbox}%
}
\renewenvironment{alertblock}[1]{%
  \pause%
  \begin{tcolorbox}[stplbase, colframe=stplred!75!black,
    colback=stplred!10!stplcard, title={#1}]%
}{%
  \end{tcolorbox}%
}
\renewenvironment{exampleblock}[1]{%
  \pause%
  \begin{tcolorbox}[stplbase, colframe=stplmuted, title={#1}]%
}{%
  \end{tcolorbox}%
}

%% ---- Code listings (Rust + Anchor) ----
\lstdefinelanguage{Rust}{
  keywords={fn,pub,struct,impl,trait,let,mut,self,Self,match,return,for,in,
            if,else,enum,use,mod,as,where,Option,Result,Some,None,Ok,Err,
            u8,u16,u32,u64,i64,u128,bool,String,Vec,Pubkey,Context,msg,vec,
            Account,Signer,Program,System,has_one,init,payer,space,seeds,bump},
  morecomment=[l]{//},
  morestring=[b]",
  sensitive=true,
}
\lstset{
  basicstyle=\ttfamily\scriptsize\color{stplwhite},
  backgroundcolor=\color{stplcard},
  keywordstyle=\color{stplred}\bfseries,
  commentstyle=\color{stplmuted},
  stringstyle=\color{stplmuted},
  showstringspaces=false,
  keepspaces=true,
  columns=fixed,
  basewidth=0.46em,
  frame=leftline,
  framerule=2pt,
  rulecolor=\color{stplred},
  xleftmargin=0.6em,
  aboveskip=0.4em,
  belowskip=0.2em,
  language=Rust,
  escapeinside={(*@}{@*)},
}

%% ---- Code reveal card (line-synced to bullets) ----
%% A block-styled card WITHOUT the block's auto-\pause, so it is visible from
%% the current overlay and its code lines can reveal in step with the itemize
%% bullets via (*@\onslide<k->@*) markers. The listing drops its own
%% background/frame (style=reveal) so the solid card shows through instead of
%% the per-line background going ragged as lines are covered.
\newtcolorbox{codeblock}[1]{stplbase, colframe=stplred, title={#1}}
\lstdefinestyle{reveal}{backgroundcolor={}, frame=none, xleftmargin=0pt,
  aboveskip=0pt, belowskip=0pt}

%% ---- Table of contents ----
\setbeamercolor{section in toc}{fg=stplwhite}
\setbeamercolor{subsection in toc}{fg=stplmuted}
\setbeamertemplate{section in toc}{%
  {\color{stplred}\textbf{\textendash}}\hspace{0.5em}\inserttocsection%
}

%% ---- Footline: minimal ----
%% Transparent bg so the background image shows through (normal text has
%% bg=stpldark, which would paint a flat band over the image's bottom glow).
\setbeamercolor{footline}{fg=stplwhite, bg=}
\setbeamertemplate{footline}{%
  \begin{beamercolorbox}[wd=\paperwidth, ht=0pt, dp=0.55cm]{footline}%
    \hspace{0.85cm}%
    \raisebox{-0.42cm}{\includegraphics[height=0.42cm]{images/stpl-logo.png}}%
    \hfill%
    \raisebox{-0.24cm}{\textcolor{stplmuted}{\tiny\insertframenumber\,/\,\inserttotalframenumber}}%
    \hspace{0.85cm}%
  \end{beamercolorbox}%
}
\setbeamertemplate{navigation symbols}{}

%% ---- Title page defaults (decks may override) ----
\date{\the\year}

%% ---- Title frame helper (uses hero background) ----
\newcommand{\maketitleframe}{%
  \stplherobg
  \begin{frame}[plain]%
    \titlepage%
  \end{frame}%
  \stpldefaultbg
}

%% ---- Section divider slides ----
\AtBeginSection[]{%
  \stplherobg
  \begin{frame}[plain]
    \begin{tikzpicture}[remember picture, overlay]
      \fill[stplred, opacity=0.07]
        ([xshift=-2.5cm]current page.east) circle (5.5cm);
      \fill[stplred]
        (current page.north west) rectangle ([xshift=5pt]current page.south west);
    \end{tikzpicture}%
    \vfill
    \hspace{0.1\paperwidth}%
    \begin{minipage}{0.8\paperwidth}
      \textcolor{stplmuted}{\small Section \thesection}\\[0.35cm]
      {\color{stplwhite}\LARGE\bfseries\insertsectionhead}\\[0.25cm]
      {\color{stplred}\rule{3cm}{1.5pt}}
    \end{minipage}%
    \vfill
  \end{frame}%
  \stpldefaultbg
}

%% ---- Title page ----
\setbeamertemplate{title page}{%
  \begin{tikzpicture}[remember picture, overlay]
    \fill[stplred]
      (current page.north west) rectangle ([xshift=5pt]current page.south west);
    \fill[stplred, opacity=0.07]
      ([xshift=-3cm,yshift=3cm]current page.north east) circle (5.5cm);
    \fill[stplred, opacity=0.04]
      ([yshift=-1.5cm]current page.south east) circle (3.5cm);
  \end{tikzpicture}%
  \vfill
  \hspace{0.08\paperwidth}%
  \begin{minipage}{0.85\paperwidth}
    \includegraphics[width=3cm]{images/stpl-logo.png}\par
    \vspace{0.8cm}
    {\usebeamerfont{title}\usebeamercolor[fg]{title}\inserttitle}\par
    \vspace{0.2cm}
    {\color{stplred}\rule{3cm}{1.5pt}}\par
    \vspace{0.3cm}
    {\usebeamerfont{subtitle}\usebeamercolor[fg]{subtitle}\insertsubtitle}\par
    \vspace{0.65cm}
    {\usebeamerfont{author}\usebeamercolor[fg]{author}\insertauthor}\par
    \vspace{0.12cm}
    {\small\textcolor{stplmuted}{\insertinstitute\enspace·\enspace\insertdate}}
  \end{minipage}%
  \vfill
}


\title{Deck Title}
\subtitle{Subtitle line}
\author{Author Name}
\institute{Organisation}

\begin{document}

\maketitleframe

\begin{frame}{Plan}
  \tableofcontents
\end{frame}

% ================================================================
\section{First Topic}

% Concept slide: bullets + exampleblock (illustration / effect).
\begin{frame}{Short Title}
  \begin{itemize}
    \item A point with its \textbf{key term} in bold
    \item Second point --- one line
    \item The problem stated first, then named: \textbf{Technical Term (TT)}
  \end{itemize}
  \vspace{0.3cm}
  \begin{exampleblock}{From everyday life}
    A concrete illustration the audience already understands.
  \end{exampleblock}
\end{frame}

% Concept slide: bullets + alertblock (problem / risk).
\begin{frame}{Second Title}
  \begin{itemize}
    \item The mechanism in three steps
    \item What the mechanism means for the user
    \item What it means for the system as a whole
  \end{itemize}
  \vspace{0.3cm}
  \begin{alertblock}{Problem}
    The short tension that pushes the narrative to the next slide.
  \end{alertblock}
\end{frame}

% ================================================================
\section{Comparisons and Formulas}

% Two columns: ONLY for direct A-vs-B opposition. No helper block needed.
\begin{frame}{Comparison}
  \begin{columns}
    \column{0.48\textwidth}
      \begin{block}{Option A}
        \begin{itemize}
          \item First trait
          \item Second trait
        \end{itemize}
      \end{block}
    \column{0.48\textwidth}
      \begin{alertblock}{Option B}
        \begin{itemize}
          \item First drawback
          \item Second drawback
        \end{itemize}
      \end{alertblock}
  \end{columns}
\end{frame}

% Math: named formula in a block, worked numbers in an exampleblock.
\begin{frame}{Formula}
  \begin{block}{Name of the formula}
    \[
      y = f(x)
    \]
    One line on what it is and what the variables mean.
  \end{block}
  \vspace{0.3cm}
  \begin{exampleblock}{Worked example}
    Real numbers substituted in: $f(2) = 4$.
  \end{exampleblock}
\end{frame}

% ================================================================
\section{Diagrams and Code}

% Static frame: <1> renders the frame on ONE overlay. Use it whenever a
% frame is a diagram or table that should land all at once.
% CAUTION: with `<1>` every bullet after the first is swallowed unless each
% is tagged \item<1->. A diagram-only frame like this one has no bullets.
\begin{frame}<1>{Static Diagram}
  \begin{center}
    \begin{tikzpicture}[
      every node/.style={font=\scriptsize\bfseries},
      box/.style={draw=dpaccent, line width=1pt, fill=dpcard, text=dpfg,
                  rounded corners=2pt, minimum width=1.4cm,
                  minimum height=0.55cm, align=center},
      arr/.style={-latex, draw=dpaccent, line width=1pt},
      lbl/.style={font=\tiny\itshape, text=dpmuted},
    ]
      \node[box] (a) {Client};
      \node[box, right=1.8cm of a] (b) {Server};
      \node[box, right=1.8cm of b] (c) {Database};
      \draw[arr] (a) -- node[lbl, above] {request} (b);
      \draw[arr] (b) -- node[lbl, above] {query} (c);
    \end{tikzpicture}
  \end{center}
\end{frame}

% Staged diagram: TikZ scopes share the bullets' overlay numbers, so the
% picture advances in lockstep with the text.
\begin{frame}{Staged Diagram}
  \begin{itemize}
    \item<1-> The starting state
    \item<2-> What the second step adds
  \end{itemize}
  \vspace{0.3cm}
  \begin{center}
    \begin{tikzpicture}[
      every node/.style={font=\scriptsize\bfseries},
      box/.style={draw=dpaccent, line width=1pt, fill=dpcard, text=dpfg,
                  rounded corners=2pt, minimum width=1.4cm,
                  minimum height=0.55cm, align=center},
      arr/.style={-latex, draw=dpaccent, line width=1pt},
      lbl/.style={font=\tiny\itshape, text=dpmuted},
    ]
      \node[box] (a) {Stage 1};
      % Declared now (invisible) so the bounding box never reflows.
      \node[box, right=1.8cm of a, visible on=<2->] (b) {Stage 2};
      \begin{scope}[visible on=<2->]
        \draw[arr] (a) -- node[lbl, above] {then} (b);
      \end{scope}
    \end{tikzpicture}
  \end{center}
\end{frame}

% Code revealed in step with the bullets. Frame MUST be [fragile].
\begin{frame}[fragile]{Code Reveal}
  \begin{itemize}
    \item<1-> The signature and what it takes
    \item<2-> The line that does the work
  \end{itemize}
  \vspace{0.3cm}
  \begin{codeblock}{example.py}
\begin{lstlisting}[style=reveal, language=Python]
(*@\onslide<1->@*)def transfer(src, dst, amount):
(*@\onslide<2->@*)    src.balance -= amount
(*@\onslide<2->@*)    dst.balance += amount
\end{lstlisting}
  \end{codeblock}
\end{frame}

% ================================================================
\section{Summary}

\begin{frame}{Summary}
  \begin{block}{The headline}
    \begin{itemize}
      \item One sentence for the first section
      \item One sentence for the second section
    \end{itemize}
  \end{block}
\end{frame}

% ================================================================
% Mandatory sign-off. Nothing goes after it.
\makeclosingframe

\end{document}
